\section{Đặc tả giao thức}
\begingroup
\setstretch{1.12}

\subsection{Cách đóng gói và đọc thông điệp}

Ứng dụng sử dụng giao thức riêng ở tầng ứng dụng, hiện tại là phiên bản 6.
Các lệnh như đăng nhập, gửi tin nhắn hay chấp nhận file được biểu diễn bằng JSON.
Trước phần JSON là một số nguyên 4 byte cho biết có bao nhiêu byte dữ liệu cần đọc:

\begin{center}
\fbox{\texttt{[Độ dài JSON: 4 byte][Nội dung JSON: N byte]}}
\end{center}

Độ dài được tính theo byte UTF-8, không tính theo số ký tự. Bốn byte độ dài được
lưu theo thứ tự byte lớn trước (big-endian). Bên nhận đọc đủ bốn byte này, lấy
ra N rồi tiếp tục đọc đủ N byte của JSON. Cách làm này cần thiết vì TCP chỉ cung
cấp một dòng byte: một thông điệp có thể đến qua nhiều lần đọc, hoặc nhiều
thông điệp có thể đến trong cùng một lần đọc. Mỗi JSON được giới hạn ở 1 MiB.

\begin{table}[H]
\centering\small
\begin{tabularx}{\linewidth}{p{3.1cm} p{2cm} >{\raggedright\arraybackslash}X}
\toprule
Trường & Kiểu dữ liệu & Ý nghĩa \\
\midrule
\texttt{opcode} & Số nguyên & Xác định nhóm chức năng, ví dụ đăng nhập, chat hoặc file. \\
\texttt{sequence} & Số nguyên & Số thứ tự gửi, tăng dần riêng ở mỗi chiều của kết nối TCP. \\
\texttt{body} & Đối tượng & Chứa loại thao tác và dữ liệu của thông điệp. \\
\texttt{body.type} & Chuỗi & Tên thao tác cụ thể trong nhóm chức năng. \\
\texttt{body.replyTo} & Số nguyên & Cho biết phản hồi đang trả lời thông điệp có số thứ tự nào. \\
\texttt{body.status} & Chuỗi & Kết quả xử lý, chẳng hạn \texttt{OK} hoặc mã lỗi. \\
\bottomrule
\end{tabularx}
\caption{Cấu trúc một thông điệp điều khiển. Hai trường cuối chỉ xuất hiện ở những loại thông điệp cần phản hồi tương ứng.}
\end{table}

Ví dụ dưới đây là yêu cầu gửi tin nhắn trong community, trước khi được mã hóa:
\noindent\begin{minipage}{\linewidth}
\begin{lstlisting}[upquote=true,columns=fullflexible,keepspaces=true]
{
  "opcode": 4,
  "sequence": 10,
  "body": {
    "type": "COMMUNITY_BROADCAST",
    "communityId": "<community-id>",
    "content": "Hello everyone"
  }
}
\end{lstlisting}
\end{minipage}

Server kiểm tra người gửi có thuộc community đó hay không. Nếu hợp lệ, server
thêm trường \texttt{sender} từ tài khoản đã đăng nhập rồi chuyển tin đến các
thành viên online khác. Server không tin tên người gửi do client tự ghi trong
nội dung yêu cầu. Với global chat, loại thông điệp là \texttt{GLOBAL\_BROADCAST};
client chỉ cần gửi \texttt{content}, vì phạm vi nhận là mọi người đang online.

\subsection{Các nhóm thông điệp}

Một \texttt{opcode} xác định nhóm chức năng, còn \texttt{type} xác định thao tác
trong nhóm đó. Ví dụ, nhóm AUTH có thao tác đăng nhập và đăng xuất.
Bảng sau tóm tắt tám nhóm được dùng trong chương trình.

\begin{table}[H]
\centering\small
\begin{tabularx}{\linewidth}{c l >{\raggedright\arraybackslash}X}
\toprule
Mã & Nhóm & Chức năng và ví dụ \\
\midrule
0 & SYSTEM & Xác nhận mã hóa, kiểm tra kết nối bằng PING/PONG và báo lỗi. \\
1 & CONNECTION & Chào kết nối bằng HELLO/HELLO\_ACK; thiết lập kết nối P2P. \\
2 & AUTH & Đăng nhập và đăng xuất bằng yêu cầu cùng phản hồi tương ứng. \\
3 & USER & Lấy danh sách online, thông báo online/offline và tìm địa chỉ peer. \\
4 & CHAT & Tin riêng MESSAGE, tin community COMMUNITY\_BROADCAST, tin chung GLOBAL\_BROADCAST và xác nhận nhận tin riêng RECEIPT. \\
5 & FILE & Đề nghị gửi, chấp nhận, truyền, kết thúc và hủy file; gồm cả file riêng và file community. \\
6 & COMMUNITY & Tạo, tham gia, rời nhóm; lấy thông tin và danh sách thành viên. \\
7 & DISCOVERY & Tìm server qua UDP bằng DISCOVERY\_REQUEST và DISCOVERY\_RESPONSE. \\
\bottomrule
\end{tabularx}
\caption{Các nhóm chức năng trong giao thức.}
\end{table}

Với tin nhắn riêng, \texttt{MESSAGE} chứa \texttt{messageId} để nhận diện tin và
\texttt{content} là nội dung. Khi nhận được tin, client gửi \texttt{RECEIPT}
kèm cùng mã tin và giá trị \texttt{DELIVERED}. Xác nhận này cho biết ứng dụng
phía nhận đã nhận tin, không có nghĩa người dùng đã đọc. Tin global và community
hiện chưa có xác nhận nhận hoặc đọc riêng cho từng người.

Các phản hồi được ghép với yêu cầu bằng \texttt{replyTo} và loại thông điệp,
giúp client không nhầm phản hồi đăng nhập với phản hồi lấy danh sách người dùng.
Những lỗi như sai mật khẩu hoặc tên community đã tồn tại được trả bằng
\texttt{status}. Riêng tên nhóm trùng sử dụng mã \texttt{NAME\_TAKEN}, để giao diện
thông báo và cho người dùng nhập tên khác.

\subsection{Thông điệp và dữ liệu truyền tập tin}

Trước khi truyền file, bên gửi thông báo các thông tin cơ bản của file.
Nhờ đó, bên nhận biết tên, kích thước và có thể lựa chọn tải về hoặc từ chối.

\begin{table}[H]
\centering\small
\begin{tabularx}{\linewidth}{p{3.8cm} >{\raggedright\arraybackslash}X}
\toprule
Trường & Ý nghĩa \\
\midrule
\texttt{transferId}/\texttt{fileId} & Mã để phân biệt các lượt truyền hoặc file được lưu trên server. \\
\texttt{sender}, \texttt{receiver} & Danh tính hai phía của file riêng. Offer community dùng người gửi và mã nhóm để xác định phạm vi nhận. \\
\texttt{fileName} & Tên file được hiển thị cho người nhận. \\
\texttt{fileSize} & Tổng số byte dữ liệu cần truyền. \\
\texttt{checksum} & Mã kiểm tra SHA-256, dùng để phát hiện file bị thay đổi hoặc nhận thiếu. \\
\bottomrule
\end{tabularx}
\caption{Các thông tin chính đi kèm một file.}
\end{table}

File riêng bắt đầu bằng \texttt{FILE\_OFFER}. Người nhận trả lời
\texttt{FILE\_DECISION} với lựa chọn \texttt{ACCEPT} hoặc \texttt{REJECT}.
Sau khi được chấp nhận, bên gửi mở socket dữ liệu riêng, gửi
\texttt{FILE\_TRANSFER\_START} và chờ \texttt{FILE\_READY} trước khi truyền.

File community sử dụng các thông điệp bắt đầu bằng \texttt{FILE\_BROADCAST\_}.
Client xin phép upload, nhận mã file cùng một mã cấp quyền tạm thời (token),
rồi mở socket riêng để tải file lên server. Khi file được lưu thành công, server
gửi \texttt{FILE\_BROADCAST\_OFFER} cho các thành viên online. Người muốn tải file gửi
\texttt{FILE\_BROADCAST\_DECISION}; server kiểm tra quyền thành viên và cấp token tải về.
Token chỉ dùng một lần và hết hạn sau 5 phút, nên không cần gửi lại mật khẩu
trên từng kết nối file.

Dữ liệu file được đọc và ghi theo khối 64 KiB, không chuyển từng khối thành JSON.
Bên nhận dựa vào \texttt{fileSize} để đọc đúng số byte, sau đó nhận
\texttt{FILE\_END} và kiểm tra SHA-256. Nếu file hợp lệ, bên nhận lưu file hoàn
chỉnh rồi gửi \texttt{FILE\_RESULT} với kết quả \texttt{OK}. Khi dữ liệu sai hoặc
thiếu, file tạm được xóa và kết quả là lỗi. File rỗng vẫn có bước kết thúc và
xác nhận; giới hạn kích thước cho mỗi file là 500 MiB.

\subsection{Mã hóa và xử lý kết nối bị gián đoạn}

Khi bắt đầu kết nối, hai phía trao đổi thông điệp HELLO để thống nhất phiên bản
và thiết lập khóa mã hóa. Client có sẵn khóa công khai tin cậy của server để
kiểm tra đúng server cần kết nối. Chương trình dùng X25519 để trao đổi khóa,
Ed25519 để kiểm tra chữ ký xác nhận danh tính và AES-256-GCM để mã hóa dữ liệu.
Các thao tác mật mã sử dụng thư viện chuẩn của Java, không tự xây dựng thuật toán.
Sau bước thiết lập ban đầu, hai bên trao đổi thông điệp xác nhận khóa trước khi
đăng nhập hoặc gửi nội dung.

Phần JSON và dữ liệu file sau đó được đặt trong các khối dữ liệu đã mã hóa.
Vì vậy, cấu trúc JSON ở trên là nội dung trước khi mã hóa, không phải văn bản
có thể đọc trực tiếp trên đường truyền. Với P2P, server ký xác nhận danh tính
và khóa công khai của peer, giúp hai client kiểm tra người đang kết nối.
UDP tìm server không được mã hóa và không thay thế bước xác thực TCP này.

Client gửi PING định kỳ 5 giây và chờ PONG tối đa 30 giây để kiểm tra kết nối.
Yêu cầu quản lý chờ phản hồi tối đa 10 giây; đề nghị nhận file P2P chờ tối đa
5 phút. Nếu hủy file riêng, bên hủy gửi \texttt{FILE\_CANCEL} và chờ
\texttt{FILE\_CANCEL\_ACK} tối đa 2 giây trước khi dọn kết nối dữ liệu.
Với file community, ứng dụng hủy bằng cách đóng socket upload hoặc download
đang dùng, không đóng kết nối quản lý.

Khi mất server, global và community không tiếp tục hoạt động, nhưng các kết nối
P2P còn sống vẫn có thể chat và truyền file riêng. Người dùng phải nhập lại mật
khẩu để kết nối lại server. Ứng dụng chưa tự gửi lại yêu cầu cũ hoặc tiếp tục
file từ vị trí đang dở. Đặc tả đầy đủ các loại thông điệp nằm trong
\path{docs/protocol.md} của mã nguồn.

\endgroup
